% =====================================================================
%  6. Estado del arte
% =====================================================================
\section{Estado del arte}
\label{sec:estado-arte}

El estado del arte se organiza en ocho líneas. Cada línea responde a una pregunta
orientadora y agrupa las fuentes que la sustentan, de modo que una fuente nueva se
incorpore a una línea existente en lugar de generar una sección adicional.

\begin{table}[H]
\centering
\small
\caption{Líneas de investigación y pregunta orientadora de cada una}
\label{tab:lineas}
\begin{tabularx}{\textwidth}{@{}c L{4.2cm} X@{}}
\toprule
\textbf{Línea} & \textbf{Ámbito} & \textbf{Pregunta orientadora} \\
\midrule
A & Calidad y curación de recursos educativos &
¿Qué significa que un recurso educativo sea de calidad y cómo se evalúa de forma
sistemática? \\
B & Problema educativo y carga docente &
¿Existe una necesidad demostrable de apoyo a la gestión de documentación académica? \\
C & Inferencia, contradicción y verificación documental &
¿Cómo se detectan contradicciones y se verifican afirmaciones contra evidencia? \\
D & Recuperación, \emph{embeddings} y RAG &
¿Cómo debe recuperarse la evidencia que alimenta al sistema? \\
D$'$ & Evaluación de sistemas RAG &
¿Con qué métricas se diagnostica un sistema RAG y qué tan fiables son? \\
E & Agentes, herramientas y seguridad &
¿Qué aporta un agente con herramientas y cómo se evalúa su comportamiento? \\
F & Confianza, factualidad y calibración &
¿Puede el sistema estimar cuándo su propia salida es confiable? \\
G & Supervisión humana, retroalimentación y trazabilidad &
¿Cómo mantener la autoridad humana y aprovechar la retroalimentación? \\
H & Metodología y evaluación con usuarios &
¿Cómo evaluar con rigor un artefacto y su interacción con personas? \\
\bottomrule
\end{tabularx}
\end{table}

\subsection{Líneas A y B: calidad, curación y carga docente}
\label{sec:arte-ab}

El contenido sustantivo de estas dos líneas se expuso en las
secciones~\ref{sec:contexto-carga} a \ref{sec:contexto-aied}, por ser el fundamento
del planteamiento del problema. Se retoman aquí los tres resultados que condicionan
el diseño de la investigación:

\begin{enumerate}
  \item La exactitud del contenido es la dimensión dominante de la calidad de un
        recurso educativo \parencite{leacock2007framework}, y los errores de los
        materiales se reproducen en la comprensión de los estudiantes.
  \item Ni la evaluación automática ni la revisión humana bastan por separado: la
        primera cubre la forma \parencite{clements2015open,tabares2017modelo}, la
        segunda subdetecta defectos \parencite{xie2018systematic}.
  \item La IA generativa redistribuye la carga docente más de lo que la reduce
        \parencite{su2026efficiency}, por lo que medir solo tiempo subestima la
        intensificación.
\end{enumerate}

A estos tres resultados se añade una asimetría de atención que la orientación de
política de \textcite{miao2021aieducation} documenta con claridad: las aplicaciones
de IA orientadas al docente han recibido mucho menos desarrollo que las orientadas al
estudiante, y las que existen se concentran en automatizar evaluación, detección de
plagio y tareas administrativas \parencite[p.~18]{miao2021aieducation}. El documento
advierte además que la eficacia aparente de muchas herramientas puede deberse más a
su novedad que a su sustancia, y que falta investigación acumulativa y replicable
sobre su impacto \parencite[p.~26]{miao2021aieducation} --- la misma observación sobre
el efecto novedad que \textcite{zawacki2019systematic} hacen desde la revisión de
literatura primaria.

\begin{advertencia}
\textcite{miao2021aieducation} es un documento de política anterior a la IA
generativa, sin evidencia empírica propia. En la matriz bibliográfica figuraba como
Holmes, Bialik y Fadel (2019); el documento efectivamente disponible y leído es este,
y la matriz debe corregirse en consecuencia.
\end{advertencia}

Queda por añadir una advertencia sobre el uso de la IA como apoyo cognitivo.
\textcite{atchley2024human} sostienen que la descarga cognitiva en herramientas de IA
tiene costos --- peor desempeño posterior sin la herramienta, confianza inapropiada
en el propio conocimiento, mayor vulnerabilidad a la información falsa --- y que las
personas no suelen percibirlos. Citan además un dato sobre fabricación de referencias
que ilustra el riesgo de aceptar salidas sin verificación.

\begin{advertencia}
\textcite{atchley2024human} es un ensayo de perspectiva sin datos propios; su
prioridad en la matriz bibliográfica debe bajarse de Núcleo a Apoyo. Se usa aquí como
articulación de un riesgo, no como evidencia de que ese riesgo se materialice en el
contexto de la curación docente.
\end{advertencia}

\subsection{Línea C: inferencia, contradicción y verificación documental}
\label{sec:arte-c}

Esta línea aporta la estructura formal de la tarea de detección y, sobre todo, la
medida de su dificultad.

\paragraph{La contradicción es la clase difícil.}
En inferencia a nivel documental sobre contratos,
\textcite{koreeda2021contractnli} reportan un $F_1$ de contradicción de 0,357 con
BERT-large frente a 0,834 en implicación, y de 0,405 con el mejor modelo evaluado. Con
evidencia provista por un oráculo, el $F_1$ de contradicción asciende a 0,62--0,66:
la recuperación de la evidencia correcta condiciona el resultado más que el
clasificador. Los autores identifican además dos fenómenos que degradan la exactitud:
la negación por excepción --- presente de forma local en el 12\,\% de los pares y de
forma no local, a veces a páginas de distancia, en el 7\,\% --- y la evidencia
discontinua, presente en el 28\,\% de los pares.

\paragraph{La clase neutral es la peor reconocida.}
\textcite{gubelmann2024capturing} muestran que los modelos entrenados con los
conjuntos dominantes, orientados a validez deductiva, tratan como neutral lo que es
inductivamente válido, con la consiguiente caída de \emph{recall} en textos
argumentativos reales. En su evaluación con silogismos,
\textbf{ningún modelo reconoció las relaciones neutrales con más del 25\,\% de
exactitud}, por debajo del 33\,\% que daría el azar; y con un desbalance realista ---
la mayoría de los pares de oraciones sin relación --- un modelo colapsó al 1\,\% de
exactitud por atribuir relaciones inexistentes. Añaden un hallazgo contraintuitivo:
los modelos más grandes no son necesariamente mejores en tareas lógicas.

\paragraph{Qué métricas de consistencia funcionan.}
\textcite{honovich2022true} metaevaluaron métricas de consistencia factual sobre once
conjuntos de datos estandarizados a etiquetas binarias. Los métodos basados en
inferencia y en generación y respuesta de preguntas resultaron los mejores --- AUC
promedio de 81,5 y 81,4 y 80,7 respectivamente --- frente a 72,2 del siguiente y un
desempeño pobre de las métricas de n-gramas. Tres advertencias de ese trabajo son
directamente aplicables aquí: todas las métricas se degradan con textos fuente de más
de 200 \emph{tokens}; promediar la consistencia por partes permite que un texto
mayormente correcto con un error pequeño reciba una puntuación alta; y entre los
ejemplos más difíciles, 35 de 80 correspondían a etiquetas humanas equivocadas.

\paragraph{La granularidad importa.}
\textcite{min2023factscore} descomponen el texto generado en hechos atómicos y
verifican cada uno contra evidencia. Su justificación es empírica: en el 40\,\% de los
casos una sola oración mezcla información respaldada y no respaldada, por lo que
evaluar por oración o por texto completo es insuficiente. Reportan además que la
precisión factual cae con entidades poco frecuentes y con los hechos mencionados más
tarde en el texto, y que su estimador automático aproxima la anotación humana con
menos de 2\,\% de error cuando dispone de recuperación de evidencia.

\subsection{Línea D: recuperación, \emph{embeddings} y generación aumentada}
\label{sec:arte-d}

\paragraph{No hay un recuperador universalmente superior.}
\textcite{karpukhin2020dense} documentan una ventaja de 9 a 19 puntos absolutos en
exactitud top-20 de la recuperación densa sobre BM25, con la excepción de conjuntos
donde la coincidencia léxica es alta. Señalan que ambos recuperan pasajes
cualitativamente distintos --- BM25 es sensible a palabras clave muy selectivas, la
recuperación densa capta variaciones léxicas y relaciones semánticas --- y que la
combinación de ambos mejora los resultados en algunos casos. El propio trabajo
fundacional de RAG lo confirma en sentido inverso: en FEVER, BM25 funcionó mejor que
el recuperador denso, posiblemente porque las afirmaciones se centran en entidades
\parencite[p.~7]{lewis2020retrieval}.

\paragraph{No hay un \emph{embedding} universal.}
\textcite{muennighoff2023mteb} concluyen que ningún modelo domina en todas las tareas
de su banco de pruebas y que existen dos familias con comportamientos distintos: los
orientados a recuperación asimétrica --- consulta corta contra documento --- y los
orientados a similitud simétrica entre textos comparables. Los modelos ajustados a
similitud textual rinden mal en recuperación y viceversa, y el desempeño cae en
idiomas no vistos durante el preentrenamiento. \textcite{reimers2019sentencebert}
aportan la base de la familia simétrica y un dato de eficiencia relevante para
detectar redundancia en un corpus: encontrar el par más similar entre diez mil
oraciones pasa de unas 65 horas con un \emph{cross-encoder} a unos 5 segundos con
\emph{embeddings} de oración.

\begin{advertencia}
Las tareas de recuperación de MTEB son exclusivamente en inglés
\parencite{muennighoff2023mteb}. La selección de un modelo de \emph{embeddings} para
un corpus en español no puede justificarse citando su posición en ese banco de
pruebas: requiere validación propia.
\end{advertencia}

\paragraph{Qué aporta y qué no aporta el patrón RAG.}
\textcite{lewis2020retrieval} documentan que evaluadores humanos juzgaron más factual
a su sistema que a la alternativa sin recuperación en el 42,7\,\% de los casos, frente
al 7,1\,\% en sentido contrario, y que el conocimiento se actualiza intercambiando el
índice documental sin reentrenar. Pero los propios autores acotan el alcance al
advertir que la fuente de conocimiento externa difícilmente será alguna vez
completamente factual y libre de sesgos \parencite[p.~10]{lewis2020retrieval}, y el
trabajo \textbf{no aborda qué ocurre cuando la fuente se contradice internamente}, que
es precisamente el caso de esta investigación.

\textcite{gao2024retrieval} sistematizan los desafíos: recuperación con baja precisión
o \emph{recall}, redundancia al recuperar información similar de varias fuentes,
incoherencia al integrar fragmentos, y el compromiso de la fragmentación --- bloques
grandes aportan contexto y ruido, bloques pequeños lo contrario. Registran además dos
elementos que este proyecto adopta: el uso de metadatos (archivo, página, fecha) para
filtrar y construir recuperación sensible al tiempo, y la observación de que la
recuperación mantiene el proceso observable al permitir que el usuario localice la
referencia original, frente a un contexto largo sin recuperación que opera como caja
negra.

\subsection{Línea D$'$: evaluación de sistemas RAG}
\label{sec:arte-dprima}

Tres marcos de evaluación compiten en este espacio y sus límites son conocidos.

\textcite{es2024ragas} proponen una evaluación sin referencia basada en descomponer la
respuesta en afirmaciones y verificarlas una a una. Reportan concordancia con
anotación humana de 0,95 en fidelidad, 0,78 en relevancia de la respuesta y 0,70 en
relevancia del contexto, y señalan que la relevancia del contexto es la dimensión más
difícil, sobre todo en contextos largos.

\textcite{saadfalcon2024ares} entrenan jueces ligeros y aplican inferencia
potenciada por predicción para acotar la estimación con intervalos de confianza:
estiman la tasa de alucinación con un error de 2,5 puntos usando 200 anotaciones, un
78\,\% menos que un enfoque basado solo en anotación humana.

\textcite{ru2024ragchecker} reportan mejor correlación con el juicio humano que los
marcos anteriores --- correlación de Pearson de 61,93 frente a 48,31 del mejor
comparador en la valoración global --- aunque todavía por debajo de la correlación
entre anotadores humanos. Sus hallazgos de diagnóstico son los más útiles para el
diseño del sistema: el recuperador importa en todas las configuraciones; un contexto
más informativo aumenta la fidelidad y reduce la alucinación; pero al aumentar el
\emph{recall} del recuperador los generadores se vuelven más sensibles al ruido
contenido en los propios fragmentos relevantes. Describen el resultado como un
trilema entre utilización del contexto, sensibilidad al ruido y fidelidad.

\begin{advertencia}
Los jueces automáticos de ARES se generalizan a cambios moderados de dominio pero
fallan ante cambios drásticos: al pasar del inglés a otros idiomas, incluido el
español, la $\tau$ de Kendall cae a 0,33 \parencite[p.~8]{saadfalcon2024ares}. Ninguno
de los tres marcos puede adoptarse para un corpus en español sin validar antes sus
jueces contra anotación humana propia.
\end{advertencia}

\subsection{Línea E: agentes, herramientas y seguridad}
\label{sec:arte-e}

\paragraph{Qué aporta el patrón de razonamiento y acción.}
\textcite{yao2023react} reportan en verificación de hechos 60,9\,\% de exactitud
frente a 56,3\,\% del razonamiento encadenado, y hasta 64,6\,\% al combinar ambos. Más
relevante que la diferencia de exactitud es el análisis de errores: en el razonamiento
sin herramientas la alucinación es el principal modo de fallo (56\,\% de los errores),
mientras que las trayectorias con acciones resultan más ancladas en evidencia. Los
modos de fallo propios del patrón --- bucles de razonamiento en el 47\,\% de los casos
analizados y búsquedas no informativas en el 23\,\% --- definen qué debe instrumentarse
y medirse.

\paragraph{Cómo se evalúa un agente.}
\textcite{mohammadi2025evaluation} proponen una evaluación holística: completar la
tarea no basta si la salida es de mala calidad, lenta, costosa, inconsistente o
insegura. Distinguen tres métodos de cálculo --- basado en código, juez automático y
evaluación humana --- con compromisos explícitos entre reproducibilidad, escalabilidad
y validez, y subrayan que, por el carácter estocástico de los agentes, la misma tarea
debe repetirse varias veces para estimar consistencia.

\paragraph{Por qué el juez automático no puede asumirse válido.}
\textcite{thakur2025judging} evaluaron jueces automáticos contra anotación humana. Solo
los modelos mayores lograron buena alineación y aun así quedaron por debajo del acuerdo
entre humanos --- el mejor, ocho puntos por debajo. Su advertencia metodológica es
precisa: el porcentaje de acuerdo engaña, porque jueces con más del 90\,\% de acuerdo
pueden desviarse más de diez puntos de la puntuación humana, por lo que debe
reportarse un coeficiente corregido por azar. Documentan además un sesgo de indulgencia
--- los jueces tienden a marcar \enquote{correcto} cuando dudan --- y casos en que
aceptan respuestas trampa.

\paragraph{La superficie de ataque.}
\textcite{zhan2024injecagent} miden la vulnerabilidad de agentes con herramientas a
instrucciones inyectadas en el contenido que procesan: 24\,\% en el escenario base y
47\,\% con refuerzo del ataque para el modelo más capaz evaluado. Dos hallazgos
condicionan el diseño de un sistema que ingiere documentos de terceros: los campos con
alta libertad de contenido son los más vulnerables, y los agentes ejecutaron acciones
dañinas incluso cuando el \emph{prompt} exigía comportamiento seguro.

\subsection{Línea F: confianza, factualidad y calibración}
\label{sec:arte-f}

\paragraph{Un puntaje no es una probabilidad.}
\textcite{guo2017calibration} establecen que las redes neuronales modernas están
sistemáticamente descalibradas y tienden a la sobreconfianza, con error de calibración
esperado típicamente entre 4\,\% y 10\,\%. Dos consecuencias importan aquí: exactitud y
calibración no se optimizan con los mismos parámetros --- un modelo puede mejorar su
exactitud empeorando sus probabilidades ---, y todo método de calibración requiere un
conjunto de validación separado con etiquetas verdaderas. El escalamiento por
temperatura, el método más simple, no cambia la predicción ni la exactitud: solo la
confianza reportada.

\paragraph{La transparencia no garantiza confianza apropiada.}
\textcite{mehrotra2024systematic} revisan sistemáticamente las intervenciones para
fomentar confianza apropiada en la interacción humano-IA. El método más común es
aumentar la transparencia --- alrededor del 52\,\% de los estudios --- mediante
explicaciones, puntajes de confianza y comunicación de incertidumbre, pero los
resultados son mixtos: las explicaciones pueden aumentar la aceptación de la
recomendación con independencia de que sea correcta, y los puntajes de confianza mal
calibrados pueden inducir confianza inapropiada. Señalan además que solo el 14\,\% de
los estudios considera el continuo completo de sobre, sub y desconfianza, y que la
dependencia se usa con frecuencia como sustituto de la confianza sin justificarlo.

\paragraph{Qué tipo de alucinación se mide.}
\textcite{huang2025survey} distinguen alucinación de factualidad de alucinación de
fidelidad y documentan que la recuperación aumentada no la elimina. Sus dos familias
de causas --- fallas de recuperación, entre ellas las fuentes con información
incorrecta o desactualizada, y cuello de botella de generación, que incluye el
contexto ruidoso, los conflictos de conocimiento y el fenómeno de pérdida de la
información intermedia --- se traducen directamente en las variables de ablación de
la Parte~V.

\subsection{Línea G: supervisión humana, retroalimentación y trazabilidad}
\label{sec:arte-g}

\paragraph{Retroalimentación sin reentrenamiento.}
\textcite{fernandes2023bridging} sistematizan las formas de integrar retroalimentación
humana en generación de lenguaje natural. Dos de sus estrategias operan en inferencia,
sin modificar los parámetros del modelo: la \textbf{memoria de retroalimentación} ---
almacenar decisiones previas y recuperar las correspondientes a entradas similares
para guiar la generación --- y el refinamiento iterativo. Su conclusión general es que
el valor reside en la retroalimentación misma más que en el método concreto para
usarla. Documentan también los riesgos de la recolección: retroalimentación de baja
calidad, contradictoria o adversarial, desacuerdo entre expertos, sesgo de anclaje y
sesgo de positividad; y advierten que con menos de mil comparaciones la mejora es
escasa.

\paragraph{Qué permite afirmar un registro de procedencia.}
El modelo de datos de procedencia del \textcite{moreau2013provdm} define entidades,
actividades y agentes y establece una restricción que este proyecto debe respetar al
formular sus afirmaciones de trazabilidad: \textbf{la derivación no se infiere
automáticamente}. Que una actividad use una entidad A y genere una entidad B no
implica que B derive de A; debe existir influencia real. El modelo incorpora además la
noción de delegación --- un agente actúa en nombre de otro que conserva la
responsabilidad ---, que formaliza con precisión la relación entre el sistema que
propone y el docente que decide.

\paragraph{Qué exige la gobernanza.}
El marco de gestión de riesgos del \textcite{nist2023airmf} sitúa validez y
fiabilidad como base de las demás características de una IA confiable y distingue
transparencia --- qué ocurrió ---, explicabilidad --- cómo se decidió --- e
interpretabilidad --- qué significa para el usuario. Sus subcategorías se traducen en
obligaciones concretas de documentación y medición, detalladas en
§\ref{sec:just-normativa}. La guía de \textcite{miao2023guidance} añade el requisito
de validación ética y pedagógica previa a la adopción institucional y la exigencia de
que el diseñador humano verifique el conocimiento factual de lo que la IA sugiera.

\subsection{Línea H: metodología y evaluación con usuarios}
\label{sec:arte-h}

\paragraph{El artefacto no es la contribución.}
\textcite{hevner2004design} formulan la distinción que gobierna este documento: lo que
separa construir un sistema de hacer investigación en ciencia del diseño es la
contribución identificable a la base de conocimiento \parencite[p.~81]{hevner2004design}.
Su tabla de métodos de evaluación --- observacionales, analíticos,
experimentales, de prueba y descriptivos --- provee el repertorio del que la Parte~V
selecciona. Advierten que cuando el artefacto forma parte de un sistema
humano-máquina se requieren teorías y métodos del comportamiento para evaluarlo, y que
la investigación en diseño es perecedera: los avances tecnológicos pueden dejarla
obsoleta con rapidez.

\paragraph{Puntos de entrada al ciclo de investigación.}
\textcite{peffers2007design} establecen que el proceso es nominalmente secuencial pero
no rígido y admiten explícitamente distintos puntos de entrada, entre ellos el
centrado en diseño y desarrollo. Advierten contra su uso como ortodoxia rígida. Es la
base sobre la que este proyecto sostiene que tener el artefacto construido antes de
formalizar la evaluación no constituye un defecto metodológico.

\paragraph{Instrumento de carga percibida.}
\textcite{hart1988development} desarrollaron una medida multidimensional de carga de
trabajo cuyas seis subescalas explican entre el 78\,\% y el 90\,\% de la varianza de la
calificación global, con fiabilidad test--retest de 0,83. Dos propiedades la hacen
adecuada para el estudio con docentes: los pesos específicos de la tarea hacen la
puntuación más sensible a las manipulaciones experimentales que una calificación
global única, y su aplicación toma menos de un minuto para las seis calificaciones.
Los pesos son además diagnósticos: identifican cuál es la fuente principal de carga.

\begin{estado}
\small
\textbf{Síntesis de la línea H.} El repertorio metodológico existe y está
consolidado. Lo que la revisión de \textcite{bond2024meta} documenta es que el campo
de la IA en educación superior lo aplica de forma irregular. La contribución
metodológica que este proyecto puede hacer no es inventar un método, sino aplicar con
rigor uno disponible a un problema donde la evidencia es escasa.
\end{estado}
